\documentclass{article}

\title{DAG vs Blockchain
Understanding the difference between DAG vs blockchain is essential for keeping up with the evolving distributed ledger landscape.}

\begin{document}
\maketitle

\section{DAG vs Blockchain: A Look at the Pluses and Minuses}
Cryptocurrency, smart contracts, and blockchain technologies revolutionized the way people see finance. Now, directed acyclic graphs (DAGs) offer a promising new path for distributed ledger technologies. It’s unlikely that DAG will replace blockchain anytime soon. Even so, it’s an exciting alternative for businesses and other users that need predictable fees and fast transaction speeds.

Understanding the difference between DAG vs blockchain is essential for keeping up with the evolving distributed ledger landscape.

\section{How does blockchain work?}
Blockchain is a distributed ledger technology that can record and store data in a publicly-viewable, immutable chain. Among their many functions, traditional blockchain networks enable users to send tokens from one wallet to another. In most cases, each new transaction is validated by global network nodes that record transactions. The nodes confirm the transaction is valid by comparing the information against the transaction history.

Blockchains use two primary consensus algorithms to approve and validate transactions. Some networks use a proof-of-stake (POS) consensus that lets network members lock their funds (which is their stake) to “double-check” the confirmed transaction before adding a new block. Some networks use a proof-of-work (POW) consensus mechanism, resulting in higher fees and energy usage.

The chunks of immutable data stored in the blockchain database are known as blocks. Blocks contain information about numerous existing transactions and can be viewed by anyone at any time. Transactions with higher fees are usually prioritized over those with lower fees if the blockchain uses the POW consensus mechanism.

Blockchain is often used for decentralized finance, but the technology has countless applications. For example, it can be used to track physical goods in each phase of the supply chain. Additionally, blockchains that support smart contracts have numerous medical use cases, such as insurance claim processing and patient data storage.

\section{How are directed acyclic graphs used in a distributed ledger?}
Traditional graphs consist of edges, or lines, connecting pairs of vertices with no defined direction. Directed acyclic graphs have vertices (which are the nodes) and edges that never form closed loops. DAGs require topological sorting that ensures nodes are visited only after prior dependencies have occurred. Still, many DAG-based ledgers can develop nodes simultaneously, resulting in fast transaction throughput and predictable fees.

Like blockchain, DAG ledgers can support smart contracts, enhancing their usefulness for many industries. DAGs also have publicly-viewable, immutable transactions and are applicable to numerous industries.

\section{DAG vs blockchain}
To accurately compare DAG and blockchain, you must examine each technology’s functionality, strengths, and weaknesses. The two methods are similar in many ways, but key differences set them apart.

\section{Reaching consensus}
DAG-based ledger nodes reach consensus significantly faster than blockchain networks. For example, Hedera uses a gossip protocol that transmits data at lightning-fast speeds. When a node receives new information, it immediately shares it with another node. The two nodes each choose another node randomly and transmit the newly learned information. The four nodes then choose four new nodes and share the information. This process continues with the number of nodes exponentially growing each time.

On the other hand, blockchains create one new block at a time, regardless of whether they use POW or POS. Still, it’s important to note that POS blockchains are typically faster than those using the POW consensus mechanism.

\section{Transaction speed}
Blockchain transaction validation requires the creation of a new block, and since they can create only one block at a time, speed is often lacking. DAG-based ledgers’ remarkable consensus mechanism results in a faster transaction speed. Since DAGs’ nodes are developed simultaneously, transaction speeds are almost always faster than blockchain transactions.

\section{Transaction costs}
DAGs are often cheaper to use than blockchains. For example, a Hedera transaction with one signature costs around $0.0001, although these fees are subject to change. The same trade on the Ethereum network could cost around $0.90 on the low end. The low fees associated with DAG-based ledgers make them a more sensible option for large corporations or anyone that does a high volume of transactions.

\section{Degree of decentralization}
Blockchains are often easier to decentralize than DAGs. Hedera’s hashgraph consensus mechanism relies on a committee of up to 39 nodes. It is run by a Council of up to 39 well-known corporations. Compared to Ethereum’s roughly 4,500 active nodes, it could be argued that Hedera is less decentralized.

\section{Scalability}
Because DAGs can process more transactions per second with lower energy and fee requirements, they are often seen as more scalable than blockchain. DAG-based ledgers are specifically more scalable than typical blockchain networks, as they don’t rely on mining or a steep increase in the number of active nodes. An increase in daily transactions can have an adverse effect on POW ledgers. Alternatively, DAG ledgers can easily handle increased transaction volume.

\section{Popularity}
As it stands, blockchain is still the most commonly used technology for digital ledgers. Hedera, IOTA, Nano, and a few other ledgers are currently the only projects using acyclic graph technology. Still, we may see that change over time. DAGs are seen by many as an ideal replacement for blockchain, thanks to their enhanced data structure, increased throughput, ease of use, speed, and low fee structure.

\section{Energy use}
DAG-based ledgers like Hedera use significantly less energy than all POW and many POS blockchains. To put it into perspective, a single Hedera transaction consumes around 0.0001 kWh, which is about 605,000 times less energy per transaction than Ethereum required before switching to POS. Like Algorand, the Hedera Contract Service is carbon negative. On the other hand, popular POW chains like Bitcoin use around 250-950 kWh per transaction. Even POS blockchains like Tezos consume around 0.04145 kwh per transaction.

According to Ethereum co-founder Vitalik Buterin, Ethereum’s switch to POS cut global energy usage by around 0.2%. If all blockchain systems were to switch to a DAG protocol, the energy savings would be even more monumental.

\section{A step forward}
DAG-based distributed ledgers are a promising step forward for digital currencies. Hedera is one of the leading DAG ledgers and has numerous real-world applications. The gossip protocol is capable of quickly passing information through multiple nodes simultaneously. Plus, users pay fees that are low and predictable. Additionally, Hedera is managed by many of the world’s leading corporations, such as Boeing, Dentons, DLA Piper, Google, and IBM.

For those looking to dip their toes into DAG-based ledgers, you’ll be happy to know the Hedera Smart Contract Service is compatible with Ethereum Virtual Machine and uses Solidity, one of the most popular Web3 programming languages. The Hedera network is already used for medical, transportation, NFT, and numerous other applications.
\end{document}